\documentclass[11pt,letterpaper]{article}

%% --------------------------
%% PACKAGES (Space Saving & Formatting)
%% --------------------------
\usepackage[margin=1.0in]{geometry} % Standard 1-inch margins
\usepackage[utf8]{inputenc}
\usepackage[T1]{fontenc}
\usepackage{mathptmx} % Times New Roman font (standard & space-efficient)
\usepackage[protrusion=true,expansion=true]{microtype} % Improves text justification to save space
\usepackage[numbers,sort&compress]{natbib} % Compresses citations [1-3]
\usepackage{graphicx}
\usepackage{url}
\usepackage{xcolor}
\usepackage{hyperref}

%% --------------------------
%% COMPRESS SECTION SPACING
%% --------------------------
\usepackage[compact]{titlesec}
\titleformat{\section}{\large\bfseries}{\thesection}{1em}{}
\titlespacing*{\section}{0pt}{10pt}{4pt} % {left}{before}{after}
\titleformat{\subsection}{\bfseries}{\thesubsection}{1em}{}
\titlespacing*{\subsection}{0pt}{8pt}{3pt}
\titleformat{\subsubsection}{\bfseries\itshape}{}{0em}{}
\titlespacing*{\subsubsection}{0pt}{6pt}{2pt}

%% --------------------------
%% BIBLIOGRAPHY FORMATTING
%% --------------------------
\setlength{\bibsep}{0pt plus 0.3ex} % Removes space between references

%% --------------------------
%% CUSTOM HEADER (Replaces \maketitle)
%% --------------------------
\newcommand{\customheader}{
    \begin{center}
        {\Large \textbf{AI-Driven Quantum Chemistry for Materials Chemistry}} \\ \vspace{2pt}
        {\large \textbf{Alastair Price, Ph.D.}} \\ 
        \vspace{1pt}
        \textit{Department of Chemistry, University of Toronto} \\
        \textit{alastairj.price@utoronto.ca \quad | \quad +1 (416) 666 3345}
    \end{center}
    \vspace{-10pt}
    \noindent\rule{\textwidth}{0.5pt} % A clean horizontal line
    \vspace{0.25pt}
}

\begin{document}

\customheader

My research program aims to overcome the long-standing challenge of accurately modelling complex molecular and solid-state systems by developing an integrated suite of computational tools that combine dispersion‐corrected density-functional theory (DFT) with cutting-edge artificial intelligence (AI) techniques. In doing so, this program will bridge theoretical insights with experimental validation, ultimately enabling high-throughput predictions of materials properties. The proposed methods will be particularly relevant for designing semiconductors, two-dimensional (2D) materials, and defect-rich solids, which play pivotal roles in fields as diverse as nuclear chemistry, minerals processing, and optoelectronics. By systematically improving both the accuracy and efficiency of DFT calculations, as well as integrating experimental observables into the computational workflow, my program lays the foundation for accelerated discovery of next-generation functional materials.

\section{Background and Objectives of Research Program}

Over the past two decades, dispersion-corrected DFT\cite{jcp-perspective,becke2014JCP,burke2012JCP} has emerged as a central tool for computational materials chemistry, enabling predictions of crystal structures \cite{moellmann2010importance, mayo2024assessment,BT7Rank}, electronic properties \cite{grimme2016dispersion}, and intermolecular interactions \cite{xdmaims}, with moderate computational cost. However, conventional DFT methods still struggle in several key areas. For instance, they often require extensive parameter sweeps \cite{gould2022poisoning,PhysRevX.15.011013,langreth1975exchange}, struggle with delocalization error \cite{wires}, and can underperform for materials containing localized electronic states \cite{yao2015efficient}, such as transition metal d-orbitals or defects, like vacancies and dopants \cite{sabino2022intrinsic}. My doctoral research contributed to addressing these issues in molecular crystals, successfully applying dispersion-corrected DFT to predict polymorphism via accurate modelling of intermolecular interactions in organic solids \cite{xdmaims,aimscsp}. Building on that foundation, my postdoctoral work introduced adaptive DFT concepts \cite{apbe0}, where machine learning (ML) algorithms dynamically optimize simulation parameters (such as the degree of exact-exchange mixing in hybrid DFT), thereby reducing computational overhead and achieving higher fidelity in gas-phase calculations for organic molecules.

A complementary development is a machine-learning framework called “amons”\cite{amons}, originally devised for molecular systems, which decomposes large structures into local fragments (amons) that capture critical bonding environments. I plan to extend this approach to solid-state amons, aiming to describe extended networks and realistic defect concentrations in materials such as oxides and layered semiconductors. This adaptation will allow for modelling far larger systems than would be feasible with standard supercell DFT. By integrating ML models trained on carefully selected amons, we will gain predictive power across broad chemical spaces, from electronically active dopants in wide-gap oxides to surface defects in mining- and catalysis-relevant minerals.

The principal objective of my proposed research is to create an end-to-end computational framework that unites adaptive DFT strategies with solid-state amons and experimental data. In the short term, my group will develop \textbf{double hybrid and adaptive machine learning for DFT}, where the generation and leverage of high-level reference data will be used to refine dispersion-corrected DFT on-the-fly, reducing reliance on costly parameter sweeps while improving accuracy for challenging materials. In parallel, my group will develop \textbf{solid-state amons for defect-rich materials}, where the amons concept is extended to complex solids, enabling accurate yet computationally tractable modelling of low-concentration defects, doped semiconductors, and 2D materials. This approach will circumvent the size limitations of conventional plane-wave supercell calculations, making realistic defect studies (e.g., 0.1\% doping) feasible. The third objective is \textbf{integration with experimental observables}, where my group will develop protocols to incorporate real-world measurements (such as lattice constants, defect formation energies, and chemical shifts) into the model training process, creating a feedback loop that tightly couples theory and experiment.

My long-term vision is to  catalyze progress in materials chemistry on multiple fronts by developing a comprehensive suite of tools that seamlessly integrate AI with dispersion-corrected solid-state DFT in a streamlined workflow. This approach aims to achieve high predictive accuracy at efficient computational costs, thereby accelerating the discovery and design of advanced materials. 
The tools my group will develop can dramatically reduce the time and expense associated with high-level, solid-state simulations, opening doors to more precise modelling of defect-related phenomena in semiconductors, catalysts, and nuclear materials. These capabilities, in turn, have immediate implications for energy technologies, environmental sustainability, and the fundamental understanding of material properties. Ultimately, the convergence of computational efficiency, experimental feedback, and advanced AI promises to transform how researchers discover and optimize new materials, paving the way for breakthroughs in electronics, photonics, energy conversion, and beyond.

Although my doctoral research established a strong track record in molecular crystals, this proposal expands into new material classes (semiconductors, 2D systems, and defect-laden solids) and introduces AI-driven techniques to meet the growing demand for predictive modelling in nuclear, mining, and advanced manufacturing contexts.
I bring a unique combination of expertise in dispersion-corrected DFT for molecular solids, algorithmic innovations in adaptive DFT from my postdoctoral work, and experience incorporating AI methods into traditional quantum chemistry. These skills directly align with the challenges identified above—namely, bridging quantum accuracy with large-scale materials design. Further, my established collaborations with experimental scientists will facilitate the acquisition of high-fidelity observables for training and validating our computational models.

\section{Proposed Research}

\subsubsection*{Theme I: Adaptive Density Functional Theory for Solids}

This theme focuses on enhancing the accuracy of dispersion-corrected DFT calculations by integrating adaptive machine-learning (ML) techniques into electronic structure workflows\cite{apbe0}. By leveraging high-level reference data, my lab would aim to bridge the gap between theoretical accuracy and practical computational expense in solid-state molecular materials modelling.  
%This theme is dedicated to enhancing the efficiency and accuracy of dispersion-corrected DFT calculations through the integration of adaptive machine-learning (ML) techniques and their incorporation into current electronic structure workflows \cite{apbe0}. Adaptive DFT for solids represents a shift in computational materials chemistry, offering routes to address critical challenges such as delocalization and basis-set superposition errors, while preserving or lowering computational costs associated with traditional methods. In adaptive DFT, ML techniques are integrated to dynamically optimize the computational parameters during simulations, allowing the method to tailor itself to the specific chemical environment of a system. The dynamic optimization enables more accurate predictions of both bonded and non-bonded interactions, which is particularly crucial for complex solid-state systems where a delicate balance of forces determines the material properties. By leveraging high-level reference data in the form of double-hybrid functionals applied to solids,  adaptive DFT will be able to refine its parameters on-the-fly, ultimately bridging the gap between theoretical accuracy and practical computational expense.  

\subsubsection*{Project I(a): Implement Improved Double Hybrids for Solids}  
Double-hybrid (DH) functionals, which combine conventional DFT with a Møller–Plesset perturbation theory (MP2) correction, offer a rigorous treatment of electron correlation \cite{d1, d2} but are currently hindered in solid-state applications by the prohibitive cost of plane-wave codes. This project aims to implement and optimize DH functionals specifically tailored for solids within the \textbf{FHI-aims} code\cite{fhiaims}. By utilizing a numerical atom-centered orbital framework and resolution-of-the-identity (RI) techniques, we can drastically reduce the computational scaling of MP2-like calculations and enable the use of optimized, smaller basis sets.

This will be the first implementation of double hybrid functionals, the use of which has been traditionally restricted isolated molecules, to periodic solids. High-fidelity results from double hybrids will be essential not only for predictive modelling of complex systems, such as extended catalytic surfaces, or molecular materials, but also for constructing the training datasets required for the adaptive methods described below. This subproject would be undertaken by one PhD student, providing training in computational chemistry, coding, and physics, ultimately integrating their work within a widely used international code base.

%This project aims to implement and optimize double-hybrid (DH) functionals \cite{dh23} specifically tailored for solid-state systems, thereby significantly enhancing the accuracy of DFT calculations for complex materials. Double-hybrid functionals combine conventional DFT with a perturbative correction from Møller–Plesset perturbation theory (MP2), offering a more rigorous treatment of electron correlation effects \cite{dh1,dh2}. However, their application to solids has been limited due to the prohibitive computational cost of traditional plane-wave–based methods. By adopting an efficient numerical atom-centred orbital framework provided by the FHI-aims code \cite{fhiaims,fhiaimshybrid}, our implementation will reduce the computational cost associated with MP2-like calculations and enable the optimized selection of smaller basis sets. This approach not only improves the intrinsic accuracy of the method but also fills a critical gap in the field, as similar implementations are not currently being pursued elsewhere.

%The optimized DH framework will serve as a robust tool for generating high-quality computational data for solids --- data that have traditionally been restricted to hybrid and generalized gradient approximation levels. The double-hybrid approach will be essential both for predictive modeling of solids and for constructing training datasets for future methodological developments. The enhanced performance will facilitate accurate simulations of large and complex solid-state systems that are currently beyond the reach of conventional methods. Examples include defect-laden materials and extended catalytic surfaces that encompass more than just the active site. This project, therefore, represents a significant methodological advance in computational materials chemistry by enabling more realistic and cost-effective modelling of solids. This subproject will be undertaken by one PhD student and will provide training within a well-established code base that is widely used, both domestically and internationally. 

\subsubsection*{Project I(b): Implementation of Adaptive DFT for Solids}  

Standard DFT functionals rely on fixed global parameters that are rarely optimal for every chemical environment. In this project, I will develop \textbf{Adaptive DFT} algorithms that dynamically optimize these parameters (e.g., the exact-exchange mixing fraction, MP2 correlation mixing fraction, basis set coefficients, etc.) on-the-fly using machine learning. This approach, which I pioneered for gas-phase calculations\cite{apbe0}, addresses critical challenges such as delocalization error and basis-set superposition error without the computational cost of higher-level methods.

We will first employ Kernel Ridge Regression (KRR) coupled with the Many-body Distribution Functional (MBDF) representation\cite{mbdf,cmbdf} to predict optimal functional parameters based on the local atomic environments. This allows the method to resolve the delicate balance between competing interactions, such as hydrogen bonding, $\pi$-stacking, and van der Waals forces, that governs solid-state properties. The algorithms will be rigorously validated against a diverse set of conventional and new benchmarks, including GMTKN55\cite{gmtkn55}, X23\cite{x23}, OMat24\cite{omat24}, and VQM24\cite{vqm24}, ensuring robustness across broad chemical space for both gas and solid phase systems. This research will serve as a training ground for 1-2 graduate students, building expertise at the intersection of quantum mechanics and machine learning/AI. 

%In this project, I will develop adaptive machine learning algorithms that optimize the computational parameters of dispersion-corrected solid-state DFT calculations. By systematically integrating high-level reference data, the ML models will be capable of dynamically adjusting key simulation parameters depending on the chemical system under study. This approach was developed during my post-doctoral work, has shown great promise for gas-phase calculations, and will ensure that calculations maintain high accuracy while operating at reduced computational cost relative to the reference levels. A particular focus will be placed on resolving the fine balance among competing intermolecular interactions, such as hydrogen bonding, $\pi$-stacking, van der Waals forces, and delocalization error, which is critical for reliable solid-state calculations. Initially, the adaptive algorithms will be validated against well-defined benchmark systems, such as a set of 55 diverse benchmarks for main-group thermochemistry, reaction barriers, and non-covalent interactions (GMTKN55) \cite{gmtkn55}, a set of 23 molecular crystal lattice energies (X23) \cite{x23,x23b}, and a set of 110 million DFT calculations 
%of inorganic bulk materials (OMat24) \cite{omat24}, before being applied to more complex cases involving semiconducting and 2D materials. 
%Implementing adaptive DFT in solids relies on an ML pipeline capable of generalizing from limited training data. To this end, we adopt the many-body distribution functional (MBDF) representation, which encodes local atomic environments using mathematical functionals of two- and three-body distributions \cite{mbdf,cmbdf}. Given the sparse data typically available for complex solids, we will employ kernel ridge regression (KRR) \cite{vapnik1999nature,rasmussen2006gaussian, GPR_deringer} to predict optimal DFT exchange-correlation parameters (e.g., the exact-exchange mixing fraction \cite{apbe0}) with minimal overhead. 
%This research will thus form a robust platform for future methodological advances and will serve as a compelling training ground for 1–2 graduate students, who will gain expertise in computational chemistry, ML approaches, and solid-state modelling.

%These projects ultimately aim to integrate state-of-the-art adaptive hybrid density functionals (aDFT) in addition to novel DH functionals into established electronic structure packages, specifically FHI-aims. The incorporation of functionals designed to address delocalization errors, and enhance the treatment of exchange-correlation effects, within  FHI-aims will improve the reliability and efficiency of DFT calculations key to future AI-based training. 
%By embedding these functionals into an efficient computational workflow, the project will facilitate accurate predictions for a broad range of systems, particularly in the solid-state regime.  
%Benchmarking against standard datasets (GMTKN55 \cite{gmtkn55}, X23 \cite{x23,x23b}, VQM24 \cite{vqm24}, and QM9 \cite{qm9}) will ensure that the newly incorporated functionals perform robustly, thereby enabling their application in complex materials such as defect-rich materials and nuclear materials. 

\subsubsection*{Theme II: Application of ML Methods to Solid-Form Defect Materials}


This theme focuses on overcoming the scaling limitations of standard supercell DFT by extending the amons (atom-in-molecules fragmentation approach) framework\cite{amons} to condensed phases. By shifting the computational focus from massive periodic materials supercells to chemically local fragments, my lab will enable the high-fidelity modelling of complex, defect-rich solids that are currently computationally intractable. 


%In this theme, we will extend the atoms-in-molecules (amons) framework \cite{amons}, originally developed for molecular systems, to solid-state materials chemistry. This approach relies on the intelligent selection of progressively larger supercells, where each configuration serves as a representative fragment that captures essential local defect chemistry and electronic structure. We will fully integrate the amons-based approach into the FHI-aims package \cite{fhiaims}, creating a cohesive workflow for high-throughput screening of defect-rich systems. Given a bulk crystal structure and defect specification (e.g., a dopant in a layered TMDC), the software will automatically generate  at set of local fragments (amons) and perform rigorous quantum-mechanical calculations on those fragments. By focusing on smaller, parallelized computations, we can afford advanced functionals, such as hybrids or double-hybrids, for near-benchmark accuracy. Each amon fragment then contributes to a shared library of defect configurations, which an ML model uses to predict properties (dopant positions, concentrations, or strain effects) in larger systems without recalculating from scratch. This modular strategy enables detailed insights into doping asymmetries, exciton recombination, and oxygen diffusion phenomena, ultimately accelerating both fundamental research and practical materials engineering.

%Many metal oxides (e.g., ZnO, SnO$_2$, NiO, Ga$_2$O$_3$, Gd$_2$O$_3$)  serve in applications ranging from transparent electronics to radiation-resistant coatings \cite{transparentcoat,nucshield}, but pose notable modelling challenges due to strong d-electron localization and complex defect chemistry. Semi-local DFT often underestimates band gaps and misplaces localized defect states, yielding qualitatively incorrect results at realistic doping levels \cite{localized2011defect}. Layered transition-metal dichalcogenides (TMDCs), like MoS$_2$ or WS$_2$, similarly involve correlated d-states and can form intricate defects such as vacancies, grain boundaries, or charge-density waves \cite{kim2022experimental}. Modelling such low-concentration defects with conventional supercell DFT can be prohibitively expensive, especially when large simulation cells are needed to avoid artificial interactions between periodic images. By combining the amons framework with efficient ML-driven workflows, we aim to tackle these computational bottlenecks, providing accurate and scalable solutions for defect-rich oxides and 2D materials.

\subsubsection*{Project II(a): Expansion of Current Methods to Create a Defect-Focused Materials Dataset}  

My lab will extend the amons concept, originally developed for molecular machine learning, to the solid state by generating a library of progressively larger supercells, where each configuration serves as a representative fragment capturing the local electronic environment around a vacancy or dopant. We will fully integrate this approach into the FHI-aims package\cite{fhiaims}, creating a cohesive workflow for automated fragment generation and high-throughput screening.

Rather than simulating a single, massive supercell to dilute defect interactions, our software will systematically decompose a bulk crystal (e.g., a doped semiconductor, or defected molecular crystal) into a set of overlapping local fragments. As these computations are smaller and highly parallelizable, we can employ advanced functionals, such as the hybrid or double-hybrid DFT, to achieve near-benchmark accuracy on each fragment. These high-fidelity calculations will feed into a shared training library, which an ML model will use to predict properties across larger length scales. This modular strategy allows for the accurate modelling of low-concentration defects (e.g., 0.1\% substitution) and extended strain fields without the exponential cost of brute-force simulations. This project will enable detailed insights into phenomena such as doping asymmetries in wide gap-oxides and exciton recombination in layered materials, as well as the effect of deletions, defects, or substitutions in a wide range of molecular materials, effectively encoding key chemical motifs while avoiding the artifacts of periodic boundary conditions.   

%We extend the amons concept to solids by generating a series of defect-containing supercells that progressively increase in size and complexity, each capturing a specific local environment around a vacancy or dopant. Taken together, these supercells form a robust library of representative amons that can be used to train ML models. This modular dataset allows us to predict properties of much larger or more complex supercells, including those with low-level doping or extended defects, without recalculating every atomic interaction. By systematically varying defect positions, concentrations, and cell dimensions, we effectively encode the key chemical motifs behind defect formation and propagation, all while avoiding the exponential cost of brute-force simulations.

%Implementation within the FHI-aims package will support efficient property evaluation and high-throughput screening of defect-rich materials. Smaller-scale computations on these amons can be parallelized, enabling the use of advanced functionals (e.g., hybrids or double-hybrids) for near-benchmark accuracy. Each amon entry feeds into a growing ML model that accurately interpolates to different dopant positions, concentrations, or strain states—crucial for tackling phenomena like doping asymmetries in wide-gap oxides or exciton formation in layered TMDCs. This approach also alleviates the common pitfall of artificially high defect concentrations in conventional supercell DFT, enabling realistic studies of, for example, 0.1\% substitutions in large crystals. Over time, a dedicated graduate student will expand the library with additional materials and defect configurations, culminating in a publicly accessible resource that accelerates both fundamental and applied research on defect-laden solids.


\subsubsection*{Project II(b): Using Amons to Study Oxygen Defects}  

We will apply the framework developed in \textbf{Project II(a)} to solve critical electronic structure challenges in metal oxides (e.g., ZnO, Ga$_2$O$_3$, UO$_2$) and transition-metal dichalcogenides (TMDCs). These materials are vital for transparent electronics and nuclear fuels but pose severe modelling challenges due to strong $d$- or $f$-electron localization \cite{localized2011defect}. Standard semi-local DFT often underestimates band gaps and misplaces defect states. Examples include failing to predict deep in-gap states of oxygen vacancies in TiO$_2$ \cite{bharti2016formation}, and inducing shallow donor levels near the conduction band edge in ZnO \cite{KIM20121711}. Similarly, in UO$_2$, each vacancy can reduce U$^{4+}$ to U$^{3+}$, altering thermal conductivity and oxidation state \cite{dorado2012first}. Such phenomena are central to semiconductor performance, catalytic reactivity, and defect evolution in nuclear fuels, but accurately modelling them remains challenging.
 
 %Oxygen vacancies in oxides critically influence electronic structure, impacting doping, conductivity, and catalytic properties. For instance, an oxygen vacancy generates Ti$^{3+}$ centers and in-gap states in TiO$_2$ \cite{bharti2016formation}, whereas it induces shallow donor levels near the conduction band edge in ZnO \cite{KIM20121711}. Similarly, in UO$_2$, each vacancy can reduce U$^{4+}$ to U$^{3+}$, altering thermal conductivity and oxidation state \cite{dorado2012first}. Such phenomena are central to semiconductor performance, catalytic reactivity, and defect evolution in nuclear fuels, but accurately modelling them remains challenging.

Conventional plane-wave DFT often underestimates oxide band gaps, struggles with self-interaction errors, and relies on large supercells to simulate dilute defects, driving up computational costs. Methods like DFT+U \cite{wang2006oxidation,aykol2014local,macke2024orbital}, hybrid functionals \cite{gruneis2014ionization, chen2014band}, or GW \cite{chen2017accuracy} can localize defect charges and recover the correct band gap, but each introduces parameter sensitivities or further raises the computational burden. The amons-based approach addresses these limitations by modelling only the local chemical environment around the vacancy.  By building up smaller supercells, we obtain compact systems amenable to higher-level methods, providing high-accuracy data on defect-induced electronic and structural perturbations, which can then be used to train predictive ML models. Machine learning models trained on these subunits can then predict defect formation energies, ionization levels, and polaron states in the extended solid, bypassing the need for massive supercell calculations. This strategy thus balances accuracy and scalability, enabling first-principles insights into oxygen vacancies under realistic doping conditions. In essence, the amons approach uses transferable local building blocks to reconstruct accurate electronic structure for defective oxides, achieving periodic-quality predictions for band gaps and defect states in complex oxides at a fraction of the computational cost.




\subsubsection*{Theme III: Bridging the Divide Between Experimental Observables and Computational Results}

Integrating experimental observations with computational modelling is crucial to enhance the realism and accuracy of materials simulations. Current ML models for material properties---including adaptive DFT schemes and the amons approach---are predominantly trained on in silico data (e.g. DFT calculations), which limits their direct applicability to real-world conditions. To truly bridge the divide, we will incorporate high-fidelity experimental data at multiple stages of model development. This strategy will calibrate ML predictions to reflect actual material behaviour and not just reproduce computational biases.

\subsubsection*{Project III: Experimental Data Validation and Use}  
The overarching goal of this theme is to incorporate high-fidelity experimental data into our computational workflows, ensuring that the ML models guiding both adaptive DFT and solid-state amons remain tethered to real-world behaviour. Initially, we will rely on published datasets—such as well-characterized crystal structures (from the ICSD \cite{levin2018nist} or COD \cite{COD,cod2}) and documented nuclear magnetic resonance (NMR) measurements—to create a reference set of material observables. These data will serve as benchmarks for calibrating the accuracy of our computational methods, revealing where predictions from both standard DFT and ML-based approaches deviate significantly from experiment. By incorporating even a limited number of experimental points into our training data, we can systematically correct biases inherited from purely theoretical calculations, thereby improving the transferability of the ML models.

Building on these established datasets, we aim to form collaborations with laboratories that specialize in advanced characterization—particularly institutions with robust X-ray diffraction (XRD) and NMR facilities—to obtain targeted measurements for materials relevant to nuclear and mining chemistry. For instance, precise XRD lattice constants can be used to refine structural predictions, while solid-state NMR chemical shifts can improve descriptions of local bonding and hydrogen bonding networks. In practice, these observables will be integrated into the ML-driven parameter-tuning process: the computational pipeline (whether it be adaptive DFT or the amons approach) will incorporate experimental discrepancies into its loss function or fit procedure, effectively “learning” how to align predicted properties with measured data. By closing this loop between theory and experiment, our models will rapidly identify and correct systematic errors, such as underestimated band gaps or mislocalized defect states, ultimately yielding more accurate simulations of complex materials under realistic conditions.


%% --------------------------
%% REFERENCES
%% --------------------------
\newpage

\bibliographystyle{unsrt}
\bibliography{example} 

\end{document}